\subsection{Exact rational projector formulas}
Only the middle axis, $h=30$, uses the fixed $L=I+J$.
Write $\mathbf1$ for the all-ones column, $H=I-J/9$, and
$L^{-1}=I-J/(h+1)$. A source triple line has conjugated projector
\[
 P'_T=\tfrac12(t_T+3\mathbf1)
       \left(t_T-\frac{10}{3(h+1)}\mathbf1\right)^t.
\]
For an original addition envelope with common core $C$, outside set $O$,
$c=|C|\in\{1,2\}$, and $n=|O|$, vectors have common value $t$ on
$C$, support in $C\cup O$, and coordinate sum $3t$. This is a
positive nondegenerate $n$-space under $H$. Put
\[
 s_0=3-c,\quad d_0=s_0^2+(c-1)n,\quad p_O=\mathbf1_O,\quad
 w=\frac{\mathbf1_C+3\mathbf1}{s_0},\quad
 z=\frac{\mathbf1_C-10\mathbf1/[3(h+1)]}{s_0}.
\]
Direct projection onto this envelope gives
\[
 P'=\Delta_O+
 \frac{s_0^2p_Oz^t+s_0^2wp_O^t+ns_0^2wz^t-(c-1)p_Op_O^t}{d_0}.
\]
Thus its correction to the diagonal mask has rank at most two.
For integer evaluation let $w_0=\mathbf1_C+3\mathbf1$,
$z_0=3(h+1)\mathbf1_C-10\mathbf1$, and $D=3(h+1)d_0$. Then
\[
 D(P'-\Delta_O)=s_0p_Oz_0^t+3(h+1)s_0w_0p_O^t
                  +nw_0z_0^t-3(h+1)(c-1)p_Op_O^t.
\]
For core one the denominator is 372 and the correction numerator is
bounded by 10538; for core two it is $93(n+1)\le2697$ and bounded
by 9844. Source lines have denominator 186 and numerator bound 332.
These bounds apply to the correction numerator, not the whole matrix.

\subsection{All physical transition cases}
Every rank-greater-than-two transition in the selected multiset has
form $M=\Delta+Z/D$, where $\Delta$ is a zero-one diagonal mask,
and fits the following bounds:
\[
\begin{array}{l|r|r|r}
\text{transition}&\operatorname{rank}Z\le r&D\le&B:=\max|Z_{ij}|\le\\\hline
\text{envelope/complement/same-core difference}&2&75516&561108\\
\text{source line to envelope}&3&5394&29316\\
\text{core two to core one}&4&10788&344978
\end{array}
\]
For core-one same-core differences the vectors
$e_j+e_i/2$ are $H$-orthonormal, so the formula remains a mask plus
a rank-two correction on the newly added coordinates. For core two,
orthogonalizing a new vector against the old outside set gives
$e_j+(\mathbf1_C-\mathbf1_O)/(|O|+1)$; these vectors have a common
offset and Gram matrix identity plus rank one, proving the same rank-two
bound. Source growth subtracts a rank-one line projector from the
rank-two envelope correction. A core-two-to-core-one edge subtracts
two rank-two corrections, with nested outside masks. Common denominators
and the displayed single-envelope bounds give the three table rows.
Rank-one and rank-two transitions may be retained as singleton calls;
zero-rank transitions disappear.

The multiset is reconstructed from the actual DAG and orientation
matching. For an addition $z$ include $\deg(z)-1$ copies of $P_{U_z}$,
one $I-P_{U_z}$, and $P_{U_z}-P_{U_y}$ for each operand $y$.
A source contributes degree-many line projectors. Side-output growth
and line cleanup retain respectively $h-1-r$ and one singleton calls;
a total contributes $P_U$ and $I_h$.
A matched continuation preserving operand $v$ from donor $u$ to later
consumer $t$ removes the occurrences
\[
 I-P_{U_u},\quad P_{U_v},\quad P_{U_t}-P_{U_v}
\]
and inserts $P_{U_t}-P_{U_u}$ before profiling. This is a change of
physical transition occurrences, not an additive matrix identity.
It is why the old rank histogram alone cannot certify these profiles.

\subsection{A bounded-minor CRT certificate}
For any square minor of $M=\Delta+Z/D$, expand by choosing mask and
correction columns. More than $r$ correction columns give zero.
Mask columns force their rows, while a remaining $j$-square correction
determinant has absolute value at most $j!B^j$. Consequently $D^r$
times every minor is an integer of absolute value at most
\[
 \mathcal B(h,r,D,B)=\sum_{j=0}^r\binom hj j!B^jD^{r-j}.
\]
At the uniform bounds in the table the three values are
\[
 275189584885776,\quad617808735437382504,\quad
 9326311597947744700460130816.
\]
Use primes $p_1=2^{61}-1$, $p_2=2^{31}-1$, and $p_3=2^{19}-1$.
All denominators are invertible modulo each. The first two bounds are
less than $p_1$, and the third is less than
\[
 p_1p_2p_3=2596143476298333157846630848266239.
\]
A rational minor in the first two cases vanishes exactly when its
$p_1$ reduction vanishes. In the third case it vanishes exactly when
all three reductions vanish: a nonzero bounded integer divisible by
their product would contradict the strict bound. Thus each northeast
corner's rational rank is the maximum of its applicable modular ranks.
This proves zero patterns as well as nonvanishing; it is not a probable
zero test.

Lower-row/lower-column elimination chooses the first remaining nonzero
row and its rightmost nonzero entry. Clearing below and then to its
left leaves the remaining submatrix unchanged by that column cleanup.
Its pivots are determined by northeast corner ranks. Consecutive pivots
whose row and column indices both increase by one give a contiguous
identity block after diagonal scaling, and hence one recursive partial
swap by the retained compiler. The supplied calculation covers 81042
frame identifiers and 184380 distinct rank-greater-than-two matrices;
1020 require the three-prime bound and their three profiles agree.
These finite counts are certificate data, distinct from the rational
exactness argument above.

The resulting width multiplicities $(c_0,\ldots,c_{30})$ are
\begin{small}\begin{align*}
[&0,732995,44182,51426,16384,25890,12459,27850,9777,17365,\\
 &7240,17013,5820,11236,4259,9880,3026,6497,1943,4817,\\
 &1082,3171,549,1880,103,1292,0,756,30,0,30].
\end{align*}\end{small}
They satisfy
\[
 \sum_{t=1}^{30}tc_t=2843940=30\cdot94740+2\cdot870.
\]
The global characteristic uses these widths with parent dimension
34560 and replication $N/v_{30}$. Diagnostic dimensions from an
intermediate profiler are not used in this characteristic.
